% GENERATED FILE. Edit canonical lessons, then run npm run generate:books.
% canonical-source-hash: fnv1a64:be13f532cd7bfc54
% canonical-lessons: ML-C542-tolil, ML-C542-udyogam, ML-C542-sthanakkayarram, ML-C542-vetanam, ML-C542-sthanamoliyuka

\chapter{Jobs and Promotions}
\label{ch:ml-jobs-and-promotions}
\hlchaptermodality{542}

\begin{chapteropening}
\textbf{By the end of this chapter:} \emph{I can say a trade, a salaried job, a promotion, pay and to step down.}

Complete the last lesson of chapter 542: I can say a trade, a salaried job, a promotion, pay and to step down.
\end{chapteropening}

\section[\texorpdfstring{\ml{തൊഴിൽ} (toḻil)}{തൊഴിൽ (toḻil)}]{\ml{തൊഴിൽ} (toḻil) — a trade; an occupation}

\label{lesson:ML-C542-tolil}



\begin{quote}
Before the new one: say the Malayalam for to go without food, then the Malayalam for to be digested.
\end{quote}

\subsection*{You'll want to know: \ml{തൊഴിൽ}}
\textbf{\ml{തൊഴിൽ}} — \emph{toḻil} — \textquotedblleft{}a trade; an occupation\textquotedblright{}.

\begin{quote}
\ml{അച്ഛന്റെ} \ml{തൊഴിൽ} \ml{കൃഷി} \ml{ആയിരുന്നു}.
\end{quote}

\begin{quote}
\emph{acchanṟe toḻil kṛṣi āyirunnu.}
\end{quote}

\begin{quote}
Father's occupation was farming.
\end{quote}

\ml{തൊഴിലാളി}, a worker, is built on it.

\begin{cousinweb}
\textbf{\ml{തൊഴിൽ}} (\emph{toḻil}) is native Dravidian: DEDR 3524 gives it as business or an occupation (DEDR 3524).
\end{cousinweb}

\begin{grammarlens}[title={What you've built}]
One more word for describing what you lived and felt.
\end{grammarlens}

\subsection*{Your turn}
\begin{itemize}
\raggedright
  \item \emph{Say it:} \emph{toḻil}
  \item \emph{Say it:} \emph{toḻil}, once more
  \item \emph{Say it:} \emph{acchanṟe toḻil kṛṣi āyirunnu.} — Father's occupation was farming.
  \item \emph{Recall:} say the Malayalam for to go without food, then the Malayalam for to be digested, then say \emph{toḻil} again
\end{itemize}

\begin{tcolorbox}[breakable,colback=teal!4,colframe=teal!35!black,title={Before you move on}]
What does \ml{തൊഴിൽ} mean? (\textquotedblleft{}A trade; an occupation\textquotedblright{}.) Say it once more.
\end{tcolorbox}
\section[\texorpdfstring{\ml{ഉദ്യോഗം} (udyōgaṁ)}{ഉദ്യോഗം (udyōgaṁ)}]{\ml{ഉദ്യോഗം} (udyōgaṁ) — a salaried job; employment}

\label{lesson:ML-C542-udyogam}



\begin{quote}
Before the new one: say the Malayalam for to be digested, then the Malayalam for a trade.
\end{quote}

\subsection*{You'll want to know: \ml{ഉദ്യോഗം}}
\textbf{\ml{ഉദ്യോഗം}} — \emph{udyōgaṁ} — \textquotedblleft{}a salaried job; employment\textquotedblright{}.

\begin{quote}
\ml{സർക്കാർ} \ml{ഉദ്യോഗം} \ml{കിട്ടാൻ} \ml{അവൻ} \ml{ഒരുപാട്} \ml{പഠിച്ചു}.
\end{quote}

\begin{quote}
\emph{sarkkār udyōgaṁ kiṭṭān avan orupāṭŭ paṭhiccu.}
\end{quote}

\begin{quote}
He studied hard to get a government job.
\end{quote}

\ml{ഉദ്യോഗസ്ഥൻ}, an official, is the one who holds an \ml{ഉദ്യോഗം}.

\begin{cousinweb}
\textbf{\ml{ഉദ്യോഗം}} (\emph{udyōgaṁ}) is borrowed from Sanskrit \emph{udyoga}, \textquotedblleft{}an undertaking, exertion\textquotedblright{} (Monier-Williams). Malayalam narrowed it to a salaried post.
\end{cousinweb}

\begin{grammarlens}[title={What you've built}]
One more word for describing what you lived and felt.
\end{grammarlens}

\subsection*{Your turn}
\begin{itemize}
\raggedright
  \item \emph{Say it:} \emph{udyōgaṁ}
  \item \emph{Say it:} \emph{udyōgaṁ}, once more
  \item \emph{Say it:} \emph{sarkkār udyōgaṁ kiṭṭān avan orupāṭŭ paṭhiccu.} — He studied hard to get a government job.
  \item \emph{Recall:} say the Malayalam for to be digested, then the Malayalam for a trade, then say \emph{udyōgaṁ} again
\end{itemize}

\begin{tcolorbox}[breakable,colback=teal!4,colframe=teal!35!black,title={Before you move on}]
What does \ml{ഉദ്യോഗം} mean? (\textquotedblleft{}A salaried job; employment\textquotedblright{}.) Say it once more.
\end{tcolorbox}
\section[\texorpdfstring{\ml{സ്ഥാനക്കയറ്റം} (sthānakkayaṟṟaṁ)}{സ്ഥാനക്കയറ്റം (sthānakkayaṟṟaṁ)}]{\ml{സ്ഥാനക്കയറ്റം} (sthānakkayaṟṟaṁ) — a promotion}

\label{lesson:ML-C542-sthanakkayarram}



\begin{quote}
Before the new one: say the Malayalam for a trade, then the Malayalam for a salaried job.
\end{quote}

\subsection*{You'll want to know: \ml{സ്ഥാനക്കയറ്റം}}
\textbf{\ml{സ്ഥാനക്കയറ്റം}} — \emph{sthānakkayaṟṟaṁ} — \textquotedblleft{}a promotion\textquotedblright{}.

\begin{quote}
\ml{നല്ല} \ml{ജോലിക്ക്} \ml{അവൾക്ക്} \ml{സ്ഥാനക്കയറ്റം} \ml{കിട്ടി}.
\end{quote}

\begin{quote}
\emph{nalla jōlikkŭ avaḷkkŭ sthānakkayaṟṟaṁ kiṭṭi.}
\end{quote}

\begin{quote}
She got a promotion for her good work.
\end{quote}

\ml{സ്ഥാനം} a position + \ml{കയറ്റം} a climb.

\begin{cousinweb}
\textbf{\ml{സ്ഥാനക്കയറ്റം}} (\emph{sthānakkayaṟṟaṁ}) is a Malayalam compound built from \ml{സ്ഥാനം}, from Sanskrit \emph{sthāna}, \textquotedblleft{}a place, a station\textquotedblright{} (CDIAL 13753), and \ml{കയറ്റം}, \textquotedblleft{}a climb\textquotedblright{}, from the native verb \ml{കയറുക} (DEDR 1255).
\end{cousinweb}

\begin{grammarlens}[title={What you've built}]
One more word for describing what you lived and felt.
\end{grammarlens}

\subsection*{Your turn}
\begin{itemize}
\raggedright
  \item \emph{Say it:} \emph{sthānakkayaṟṟaṁ}
  \item \emph{Say it:} \emph{sthānakkayaṟṟaṁ}, once more
  \item \emph{Say it:} \emph{nalla jōlikkŭ avaḷkkŭ sthānakkayaṟṟaṁ kiṭṭi.} — She got a promotion for her good work.
  \item \emph{Recall:} say the Malayalam for a trade, then the Malayalam for a salaried job, then say \emph{sthānakkayaṟṟaṁ} again
\end{itemize}

\begin{tcolorbox}[breakable,colback=teal!4,colframe=teal!35!black,title={Before you move on}]
What does \ml{സ്ഥാനക്കയറ്റം} mean? (\textquotedblleft{}A promotion\textquotedblright{}.) Say it once more.
\end{tcolorbox}
\section[\texorpdfstring{\ml{വേതനം} (vētanaṁ)}{വേതനം (vētanaṁ)}]{\ml{വേതനം} (vētanaṁ) — pay; wages}

\label{lesson:ML-C542-vetanam}



\begin{quote}
Before the new one: say the Malayalam for a salaried job, then the Malayalam for a promotion.
\end{quote}

\subsection*{You'll want to know: \ml{വേതനം}}
\textbf{\ml{വേതനം}} — \emph{vētanaṁ} — \textquotedblleft{}pay; wages\textquotedblright{}.

\begin{quote}
\ml{തൊഴിലാളികൾ} \ml{കൂടുതൽ} \ml{വേതനം} \ml{ചോദിച്ചു}.
\end{quote}

\begin{quote}
\emph{toḻilāḷikaḷ kūṭutal vētanaṁ cōdiccu.}
\end{quote}

\begin{quote}
The workers asked for more pay.
\end{quote}

Formal, in contracts and the news; \ml{ശമ്പളം} and \ml{കൂലി} are the everyday words.

\begin{cousinweb}
\textbf{\ml{വേതനം}} (\emph{vētanaṁ}) is borrowed from Sanskrit \emph{vetana}, \textquotedblleft{}wages, hire, salary\textquotedblright{} (Monier-Williams).
\end{cousinweb}

\begin{grammarlens}[title={What you've built}]
One more word for describing what you lived and felt.
\end{grammarlens}

\subsection*{Your turn}
\begin{itemize}
\raggedright
  \item \emph{Say it:} \emph{vētanaṁ}
  \item \emph{Say it:} \emph{vētanaṁ}, once more
  \item \emph{Say it:} \emph{toḻilāḷikaḷ kūṭutal vētanaṁ cōdiccu.} — The workers asked for more pay.
  \item \emph{Recall:} say the Malayalam for a salaried job, then the Malayalam for a promotion, then say \emph{vētanaṁ} again
\end{itemize}

\begin{tcolorbox}[breakable,colback=teal!4,colframe=teal!35!black,title={Before you move on}]
What does \ml{വേതനം} mean? (\textquotedblleft{}Pay; wages\textquotedblright{}.) Say it once more.
\end{tcolorbox}
\section[\texorpdfstring{\ml{സ്ഥാനമൊഴിയുക} (sthānamoḻiyuka)}{സ്ഥാനമൊഴിയുക (sthānamoḻiyuka)}]{\ml{സ്ഥാനമൊഴിയുക} (sthānamoḻiyuka) — to step down}

\label{lesson:ML-C542-sthanamoliyuka}



\begin{quote}
Before the new one: say the Malayalam for a promotion, then the Malayalam for pay.
\end{quote}

\subsection*{You'll want to know: \ml{സ്ഥാനമൊഴിയുക}}
\textbf{\ml{സ്ഥാനമൊഴിയുക}} — \emph{sthānamoḻiyuka} — \textquotedblleft{}to step down\textquotedblright{}.

\begin{quote}
\ml{അറുപത്} \ml{വയസ്സായപ്പോൾ} \ml{മാനേജർ} \ml{സ്ഥാനമൊഴിഞ്ഞു}.
\end{quote}

\begin{quote}
\emph{aṟupatŭ vayassāyappōḷ mānējar sthānamoḻiññu.}
\end{quote}

\begin{quote}
At sixty the manager stepped down.
\end{quote}

Past: \ml{സ്ഥാനമൊഴിഞ്ഞു}, \emph{sthānamoḻiññu}.

\ml{സ്ഥാനം} a post + \ml{ഒഴിയുക} to leave, to become empty.

\begin{cousinweb}
\textbf{\ml{സ്ഥാനമൊഴിയുക}} (\emph{sthānamoḻiyuka}) is a Malayalam compound built from \ml{സ്ഥാനം}, from Sanskrit \emph{sthāna}, \textquotedblleft{}a place, a station\textquotedblright{} (CDIAL 13753), and \ml{ഒഴിയുക}, \textquotedblleft{}to leave off, to be empty\textquotedblright{}, which is native (DEDR 1009).
\end{cousinweb}

\begin{grammarlens}[title={What you've built}]
One more word for describing what you lived and felt.
\end{grammarlens}

\subsection*{Your turn}
\begin{itemize}
\raggedright
  \item \emph{Say it:} \emph{sthānamoḻiyuka}
  \item \emph{Say it:} \emph{sthānamoḻiyuka}, once more
  \item \emph{Say it:} \emph{aṟupatŭ vayassāyappōḷ mānējar sthānamoḻiññu.} — At sixty the manager stepped down.
  \item \emph{Recall:} say the Malayalam for a promotion, then the Malayalam for pay, then say \emph{sthānamoḻiyuka} again
\end{itemize}

\begin{tcolorbox}[breakable,colback=teal!4,colframe=teal!35!black,title={Before you move on}]
What does \ml{സ്ഥാനമൊഴിയുക} mean? (\textquotedblleft{}To step down\textquotedblright{}.) Say it once more.
\end{tcolorbox}
